% !TeX root = ../../main-german.tex
\chapter{Nutzung von KI-Tools}\label{sec:ai-tools}
%%% ===============================================================================

% Hier wird dokumentiert, inwiefern KI-Tools zur Erstellung wesentlicher Teile dieser Arbeit verwendet wurden.
% Die Erklärung am Ende dieser Arbeit fordert diesen Anhang.
% Siehe die Formvorgaben für Abschlussarbeiten des Fachbereichs Informatik der Universität Stuttgart: https://www.f05.uni-stuttgart.de/informatik/dokumente/Formulare/Formvorgaben-Abschlussarbeiten.pdf
% Die Struktur ist mit der Prüferin oder dem Prüfer abzustimmen.
% Nur übernommene KI-Outputs (z.B. Text, Abbildungen) müssen aufgeführt werden; die Nutzung von KI zur Recherche oder Rechtschreibprüfung nicht.
% Wurden keine KI-Outputs als wesentliche Teile übernommen, ist dies statt der Tabelle zu erklären.

\begin{tabular}{lll}
  \toprule
  KI-Tool          & Genutzt für                    & Genutzt in                \\
  \midrule
  ChatGPT5         & Erstellung von Abbildungen     & [Abbildungsreferenzen ergänzen] \\
  Google AI search & Pseudocode für Algorithmus XYZ & [Algorithmusreferenz ergänzen] \\
  DeepL Write      & Umformulierung der Einleitung  & \zcref{sec:introduction}  \\
  \bottomrule
\end{tabular}

